\documentclass[11pt,a4paper]{article}

% --------------------------------------------------
% Packages
% --------------------------------------------------

\usepackage[utf8]{inputenc}
\usepackage[T1]{fontenc}
\usepackage{lmodern}

\usepackage{amsmath}
\usepackage{amssymb}
\usepackage{amsthm}
\usepackage{mathtools}
\usepackage{bm}

\usepackage{geometry}
\usepackage{booktabs}
\usepackage{array}
\usepackage{enumitem}
\usepackage{hyperref}

\geometry{
    margin=1in
}

\hypersetup{
    colorlinks=true,
    linkcolor=blue,
    urlcolor=blue,
    citecolor=blue
}

% --------------------------------------------------
% Custom Commands
% --------------------------------------------------

\newcommand{\Cl}{\mathrm{Cl}}
\newcommand{\R}{\mathbb{R}}
\newcommand{\C}{\mathbb{C}}
\newcommand{\T}{\mathbb{T}}
\newcommand{\Spin}{\operatorname{Spin}}
\newcommand{\so}{\mathfrak{so}}
\newcommand{\G}{\mathcal{G}}
\newcommand{\D}{\mathcal{D}}
\newcommand{\F}{\mathcal{F}}
\newcommand{\W}{\mathcal{W}}
\newcommand{\Om}{\Omega}
\newcommand{\Ric}{\mathrm{Ric}}

\newcommand{\nzero}{\mathbf{n}_o}
\newcommand{\ninf}{\mathbf{n}_\infty}
\newcommand{\e}{\mathbf{e}}

\newcommand{\wedgeprod}{\mathbin{\wedge}}
\newcommand{\cross}{\mathbin{\times}}

\newcommand{\rev}[1]{\widetilde{#1}}
\newcommand{\grade}[2]{\left\langle #1\right\rangle_{#2}}

\newcommand{\comm}[2]{\left[#1,#2\right]}

% --------------------------------------------------
% Document
% --------------------------------------------------

\title{
    \textbf{Conformal Spacetime Algebra Gauge Theory of Gravity}\\
    \large A Versor-Based Extension of Gauge Theory Gravity
}

\author{}
\date{\today}

\begin{document}

\maketitle

\begin{abstract}

Gauge Theory Gravity (GTG) is formulated within the real
Conformal Spacetime Algebra (CSTA), \(\Cl(2,4;\R)\), extending the
Spacetime Algebra (STA) \(\Cl(1,3;\R)\) of conventional GTG to a
six-dimensional conformal algebra while retaining its gauge-theoretic
structure.

The gravitational gauge connection is a bivector-valued field
\(\Om(a,x)\in\Cl^2(2,4;\R)\), whose fifteen components generate
infinitesimal conformal transformations, and finite gauge
transformations and parallel transport are CSTA versors
\(V(x)\in\Spin(2,4)\). Lorentz transformations, translations,
dilations and special conformal transformations are thereby
represented in a single real Clifford algebra.

The translation sector holds the coframe
and the Lorentz sector its spin connection; GTG's Kerr solution and
its gauge freedom take exactly this form. The normal conditions fix
the special-conformal sector to the Schouten tensor, and the
quadratic curvature density is then Weyl squared. A Weyl rescaling is
a \(\Spin(2,4)\) gauge transformation, so the dilation and
special-conformal sectors carry a choice of scale. With the scale
tractor (the CSTA form of Penrose's infinity twistor) as compensator,
Einstein's equations with a cosmological constant follow both in
second-order form and, varying the full connection alone, in
first-order form, whereas the action quadratic in curvature gives
Weyl gravity, and a Weyl-squared admixture adds Stelle's massive
spin-2 ghost. Finally, the even subalgebra is \(M(4,\C)\) with the
pseudoscalar as imaginary unit, twistor space is its minimal left
ideal, and the normal conformal connection acts on it as Penrose's
local twistor transport.

\end{abstract}

\tableofcontents

\newpage

% --------------------------------------------------
\section{Overview}
% --------------------------------------------------

The formulation extends the algebraic setting of conventional Gauge
Theory Gravity from STA to CSTA. Its central features are:

\begin{enumerate}
    \item A real Clifford algebra,
    \(
        \G=\Cl(2,4;\R).
    \)

    \item A bivector-valued gravitational gauge connection,
    \[
        \boxed{
            \Om(a,x)\in\G^2.
        }
    \]

    \item Versor-valued finite gauge transformations,
    \[
        \boxed{
            V(x)\in\Spin(2,4).
        }
    \]

    \item A gauge-covariant derivative constructed using the
    geometric product.

    \item A bivector-valued curvature describing the field strength
    of the conformal connection.
\end{enumerate}

The distinction between the gauge connection and the finite gauge
transformations is fundamental. The bivector-valued connection
generates infinitesimal transformations, whereas versors represent
finite transformations and parallel transport.

% --------------------------------------------------
\section{Algebraic Foundation}
\label{sec:algebra}
% --------------------------------------------------

\subsection{Conformal Spacetime Algebra}

The underlying algebra is
\[
    \G=\Cl(2,4;\R),
\]
generated by six orthogonal basis vectors
\(
    \{\e_t,\e_x,\e_y,\e_z,\e_+,\e_-\}
\)
satisfying
\[
    \e_A\e_B+\e_B\e_A=2\eta_{AB},
    \qquad
    \eta_{AB}=\operatorname{diag}(+1,-1,-1,-1,+1,-1).
\]
Consequently
\[
    \e_t^2=+1,\qquad
    \e_x^2=\e_y^2=\e_z^2=-1,\qquad
    \e_+^2=+1,\qquad
    \e_-^2=-1.
\]
The algebra has dimension \(2^6=64\), and its bivector subspace has
dimension \(\binom{6}{2}=15\). These fifteen bivectors generate the
Lie algebra
\[
    \mathfrak{spin}(2,4)\cong\so(2,4),
\]
and the corresponding spin group is a double cover of the identity
component of the four-dimensional conformal group.

\subsection{Products and Notation}

For multivectors \(X,Y\),
\[
    [X,Y]=XY-YX,
    \qquad
    X\cross Y=\tfrac12[X,Y]
\]
denote the commutator and the commutator product. For a bivector \(B\)
and a vector \(v\), \(B\cross v=B\cdot v\), and the commutator product
of two bivectors is again a bivector. \(\rev{X}\) is the reverse of
\(X\) and \(\grade{X}{k}\) its grade-\(k\) part;
\(\grade{X}{0,6}\) is the sum of its scalar and pseudoscalar parts.

\subsection{Index Conventions}

Two distinct index sets are required.

Spacetime indices,
\[
    \mu,\nu,\ldots\ \text{(coordinate)},\qquad
    a,b,c,\ldots\ \text{(frame)},\qquad
    \text{both over }\{t,x,y,z\}.
\]
Conformal-space indices,
\[
    A,B\in\{t,x,y,z,+,-\}.
\]
The spacetime indices identify directions in the four-dimensional
spacetime on which the gauge connection is evaluated; the
conformal-space indices identify the internal bivector generators of
the six-dimensional gauge algebra. For example, \(\omega^{AB}(a,x)\)
carries two conformal-space indices and one spacetime direction
argument.

\subsection{The Pseudoscalar}

The pseudoscalar
\[
    I=\e_t\e_x\e_y\e_z\e_+\e_-
\]
satisfies \(I^2=-1\), commutes with every even multivector, and
anticommutes with every vector. It is the imaginary unit of the
twistor construction of Section~\ref{sec:twistors}.

% --------------------------------------------------
\section{Null Basis and Conformal Generators}
% --------------------------------------------------

Although the six-dimensional basis above is orthogonal, conformal
transformations are most naturally expressed in a null basis. Define
\[
    \boxed{
        \nzero=\frac{1}{\sqrt{2}}\left(\e_++\e_-\right),
        \qquad
        \ninf=\frac{1}{\sqrt{2}}\left(\e_+-\e_-\right).
    }
\]
These satisfy
\[
    \nzero^2=0,\qquad
    \ninf^2=0,\qquad
    \nzero\cdot\ninf=1.
\]
The resulting basis consists of four spacetime basis vectors and two
mutually null conformal basis vectors, and the bivector space
decomposes into four sectors.

\subsection{The Four Generator Sectors}

\begin{center}
\begin{tabular}{llcll}
    \toprule
    Sector & Generators & Number & Gauge field & Role\\
    \midrule
    Lorentz & \(M_{ab}=\e_a\wedgeprod\e_b\) & 6 & \(\omega^{ab}(a,x)\)
        & rotations and boosts\\
    Translation & \(P_a=\e_a\wedgeprod\ninf\) & 4 & \(e^a(a,x)\)
        & translations\\
    Special conformal & \(K_a=\e_a\wedgeprod\nzero\) & 4 & \(f^a(a,x)\)
        & special conformal\\
    Dilation & \(D=\nzero\wedgeprod\ninf\) & 1 & \(b(a,x)\)
        & scale transformations\\
    \bottomrule
\end{tabular}
\end{center}

The conformal bivector space therefore decomposes as
\[
    \so(2,4)
    =
    \so(1,3)\oplus\mathfrak{p}\oplus\mathfrak{k}\oplus\mathfrak{d},
    \qquad
    15=6+4+4+1.
\]

\subsection{Structure Relations}

The sectors close as the conformal algebra with the grading
\(\mathfrak{k}\oplus(\so(1,3)\oplus\mathfrak{d})\oplus\mathfrak{p}\):
\begin{gather*}
    [M,M]\subset M,\qquad
    [M,P]\subset P,\qquad
    [M,K]\subset K,\qquad
    [M,D]=0,\\
    [P,P]=0,\qquad
    [K,K]=0,\qquad
    [D,P]\propto P,\qquad
    [D,K]\propto K,\qquad
    [P,K]\subset M\oplus D,
\end{gather*}
with \(D\) acting on \(P\) and \(K\) with opposite signs. The Jacobi
identity holds. A dual basis follows from the reflection
\(G\mapsto nGn\), which exchanges \(\nzero\) and \(\ninf\) up to sign:
normalizing \(nG_in\) gives reciprocal generators \(G^i\) with
\(\grade{G^iG_j}{0}=\delta^i_j\), so the coordinates of a bivector
\(B\) are \(c^i=\grade{G^iB}{0}\).

\subsection{Sector Projectors}

Every bivector decomposes uniquely as
\[
    B=B_M+u\wedgeprod\ninf+v\wedgeprod\nzero+\beta\,\nzero\wedgeprod\ninf,
\]
with \(u,v\) spacetime vectors and \(B_M\) a Lorentz bivector.
Contracting with the null vectors gives
\[
    \nzero\cdot B=-u-\beta\,\nzero,\qquad
    \ninf\cdot B=-v+\beta\,\ninf,\qquad
    \beta=-\ninf\cdot(\nzero\cdot B),
\]
so the sector projectors are
\begin{align*}
    \Pi_D(B)&=-(\nzero\wedgeprod\ninf)\,
        \bigl[\ninf\cdot(\nzero\cdot B)\bigr],\\
    \Pi_P(B)&=-(\nzero\cdot B)\wedgeprod\ninf-\Pi_D(B),\\
    \Pi_K(B)&=-(\ninf\cdot B)\wedgeprod\nzero-\Pi_D(B),\\
    \Pi_M(B)&=B-\Pi_P(B)-\Pi_K(B)-\Pi_D(B).
\end{align*}
Being linear maps on bivectors, they apply equally to any
bivector-valued map, in particular to the connection and the
curvature.

\subsection{Translators and Conformal Points}

The translation generators move the origin. For an event \(x\), the
translator
\[
    T_x=\exp\left(\tfrac12\,x\wedgeprod\ninf\right)
    =1+\tfrac12\,x\,\ninf
\]
is exact at first order because \((x\wedgeprod\ninf)^2=0\). The
event's conformal point is the origin carried there,
\[
    \boxed{
        X=T_x\,\nzero\,\rev{T}_x,
    }
\]
a null vector of weight \(X\cdot\ninf=1\) from which \(x\) can be read
back.

\subsection{Action of the Generators on Conformal Points}

Acting on conformal points, \(X\mapsto VX\rev V\), the versors
\(V=\exp(cG)\) of the four sectors give:

\begin{center}
\begin{tabular}{ll}
    \toprule
    Versor & Action on events\\
    \midrule
    \(\exp(c\,P_a)\) & \(x\mapsto x+2c\,\e_a\), with \(\ninf\) fixed:
        translation\\
    \(\exp(\lambda D)\) & \(x\mapsto e^{-2\lambda}x\): dilation\\
    \(\exp(c\,K_a)\) & fixes \(\nzero\), preserves the null cone:
        special conformal\\
    \(\exp(c\,M_{ab})\) & preserves \(x^2\): Lorentz\\
    \bottomrule
\end{tabular}
\end{center}

The factors \(2\) and \(-2\) arise from the two-sided sandwich: a
versor \(\exp(\tfrac12\theta G)\) moves by the parameter \(\theta\).

% --------------------------------------------------
\section{Bivector-Valued Gauge Connection}
% --------------------------------------------------

The fundamental gauge field is
\[
    \boxed{
        \Om(a,x)
        =
        \frac{1}{2}
        \omega^{AB}(a,x)\,
        \e_A\wedgeprod\e_B
        \in\G^2,
    }
\]
a bivector-valued function, linear in the spacetime direction \(a\)
and depending on position \(x\). In the null basis,
\[
\boxed{
    \Om(a,x)
    =
    \frac{1}{2}\omega^{ab}(a,x)\,\e_a\wedgeprod\e_b
    +e^a(a,x)\,\e_a\wedgeprod\ninf
    +f^a(a,x)\,\e_a\wedgeprod\nzero
    +b(a,x)\,\nzero\wedgeprod\ninf .
}
\]
Equivalently, with \(e(a)=e^a(a)\,\e_a\) and \(f(a)=f^a(a)\,\e_a\)
spacetime vectors,
\[
    \Om=\Om_M+e\wedgeprod\ninf+f\wedgeprod\nzero
       +b\,\nzero\wedgeprod\ninf ,
\]
and each sector is read off by the projectors of the previous
section. The connection contains four types of gauge fields:

\begin{center}
\begin{tabular}{lll}
    \toprule
    Gauge sector & Field & Independent components per direction\\
    \midrule
    Lorentz & \(\omega^{ab}\) & 6\\
    Translation & \(e^a\) & 4\\
    Special conformal & \(f^a\) & 4\\
    Dilation & \(b\) & 1\\
    \midrule
    Total & & 15\\
    \bottomrule
\end{tabular}
\end{center}

Over the four spacetime directions the connection has \(60\)
components. \(\Om(a,x)\) is the infinitesimal gauge connection, not a
finite gauge transformation.

% --------------------------------------------------
\section{Versor-Valued Gauge Transformations}
% --------------------------------------------------

Finite gauge transformations are represented by CSTA versors,
\(V(x)\in\Spin(2,4)\), normalized so that
\[
    \boxed{
        V\rev{V}=1,
    }
    \qquad
    V^{-1}=\rev{V}.
\]
Every versor of the form \(V=\exp(B)\) with \(B\in\G^2\) satisfies
this, and \(V\) and \(-V\) give the same adjoint action, which
preserves grade 2 and the inner product of bivectors.

A spinor-like field transforms as
\[
    \Psi'(x)=V(x)\Psi(x),
\]
and geometric objects transforming by conjugation as
\[
    X'(x)=V(x)X(x)\rev{V}(x).
\]

\subsection{Gauge Transformation of the Connection}

The connection transforms according to
\[
    \boxed{
    \Om'(a,x)
    =
    V\Om(a,x)\rev{V}
    -
    2(a\cdot\nabla V)\rev{V}.
    }
\]
The first term is the adjoint action of the versor on the existing
connection; the second accounts for the position dependence of the
transformation. The right-hand side is again a bivector: differentiating
\(V\rev V=1\) shows that \((a\cdot\nabla V)\rev V\) is
reverse-antisymmetric, so it has only grades \(2\) and \(6\), and the
grade-6 part vanishes for versor fields \(V=\exp B(x)\).

Starting from \(\Om=0\), the transformation gives the pure-gauge
connection
\[
    \Om_{\mathrm{pure}}(a)=-2(a\cdot\nabla V)\rev{V}.
\]

% --------------------------------------------------
\section{Covariant Derivative}
% --------------------------------------------------

The covariant derivative of a spinor-like field is
\[
    \boxed{
        \D_a\Psi
        =
        a\cdot\nabla\Psi
        +
        \frac{1}{2}\Om(a,x)\Psi,
    }
\]
and that of a field transforming by conjugation
\[
    \D_aX=a\cdot\nabla X+\Om(a,x)\cross X.
\]
The factor of \(\tfrac12\) is the spin representation of bivector
generators. Both derivatives are gauge covariant,
\[
    \boxed{
        \D'_a\Psi'
        =
        V\D_a\Psi,
    }
    \qquad
    \D'_aX'=V(\D_aX)\rev V,
\]
which follows from \(\rev VV=1\). Under a pure-gauge connection,
\(V\Psi_0\) with constant \(\Psi_0\) is covariantly constant.

% --------------------------------------------------
\section{Curvature}
% --------------------------------------------------

For constant (commuting) directions \(a,b\), the curvature is
\[
\boxed{
    \F(a,b)
    =
    a\cdot\nabla\Om(b)
    -
    b\cdot\nabla\Om(a)
    +
    \frac{1}{2}[\Om(a),\Om(b)].
}
\]
For general vector fields the term \(-\Om([a,b])\) must be added. In
form language,
\[
    \F=d\Om+\tfrac12\,\Om\wedge_{\mathrm{Lie}}\Om .
\]
The curvature is bivector-valued, antisymmetric in \(a,b\), and
covariant,
\[
    \F'(a,b)=V\F(a,b)\rev{V}.
\]
It vanishes for every pure-gauge connection, and not for a generic one.

\subsection{Relation to the Covariant Derivative}

The factor \(\tfrac12\) in front of the commutator is fixed by
\[
    \boxed{
        [\D_a,\D_b]\Psi
        =
        \frac{1}{2}\F(a,b)\Psi
    }
\]
for constant commuting directions.

\subsection{Bianchi Identity}

With the covariant derivative of conjugation-type fields,
\[
    \sum_{\mathrm{cyclic}\,(l,a,b)}
    \Bigl(
        \partial_l\F(a,b)+\Om(l)\cross\F(a,b)
    \Bigr)=0 .
\]

\subsection{Curvature Decomposition}

The projectors split the curvature into four sectors,
\[
    \F=\Pi_M\F+\Pi_P\F+\Pi_K\F+\Pi_D\F .
\]
For a connection whose translation sector is a coframe
(Section~\ref{sec:normal}) they have a direct geometric content:

\begin{center}
\begin{tabular}{ll}
    \toprule
    Sector & Content\\
    \midrule
    \(\Pi_P\F\) & torsion; zero for the Levi-Civita connection\\
    \(\Pi_M\F\) & Riemann curvature minus terms in \(e\wedgeprod f\)\\
    \(\Pi_K\F\) & the exterior covariant derivative of \(f\)
        (Cotton type)\\
    \(\Pi_D\F\) & the antisymmetric part of \(e\cdot f\)\\
    \bottomrule
\end{tabular}
\end{center}

% --------------------------------------------------
\section{Versor-Valued Parallel Transport}
\label{sec:transport}
% --------------------------------------------------

The bivector-valued connection generates finite gauge transformations
through parallel transport. Transport along a path \(\gamma(s)\)
solves
\[
    \frac{d\Psi}{ds}=-\frac12\Om(\dot\gamma)\Psi,
    \qquad
    \Psi(s)=T(s)\Psi(0),
\]
so that
\[
    \boxed{
        T(x_1,x_2)
        =
        \mathcal{P}
        \exp\left(
            -\frac{1}{2}
            \int_{\gamma}\Om
        \right)
        =
        \lim_{N\to\infty}\prod_{k=N}^{1}
        \exp\bigl(-\tfrac12\Om(\dot\gamma_k)\,\Delta s\bigr),
    }
\]
with later steps multiplying on the left. Each factor is the
exponential of a bivector, so the transport is a versor,
\[
    T\rev{T}=1 .
\]
This gives the relationship between the infinitesimal and finite
descriptions:
\[
    \boxed{
        \text{Bivector connection}
        \longrightarrow
        \text{Versor-valued parallel transport}.
    }
\]

\subsection{Holonomy}

Around the small rectangle
\(x_0\to+\epsilon\partial_\mu\to+\epsilon\partial_\nu\to
-\epsilon\partial_\mu\to-\epsilon\partial_\nu\),
\[
    \boxed{
        T_{\mathrm{loop}}
        =
        \exp\left(-\tfrac12\,\epsilon^2\F(\partial_\mu,\partial_\nu)\right)
        +O(\epsilon^3).
    }
\]

\subsection{Pure Gauge and General Connections}

For a pure-gauge connection, transport is
\(T(x_1,x_2)=V(x_2)V(x_1)^{-1}\) and every holonomy is trivial. A
general connection cannot be reconstructed from a single globally
defined versor field: the connections with nonzero curvature are
exactly those with nontrivial holonomy, and they must be allowed to
exist independently of a versor potential, or be described by local
bundle sections and transition functions. This distinction is
important if versors are to become fundamental dynamical variables.

% --------------------------------------------------
\section{Relationship to Conventional Gauge Theory Gravity}
\label{sec:gtg}
% --------------------------------------------------

Conventional GTG is formulated within \(\Cl(1,3;\R)\). Its
gravitational gauge structure consists of a displacement (position)
gauge field \(h(a)\) and a bivector-valued rotation gauge field
\(\Om_{\mathrm{STA}}(a,x)\in\Cl^2(1,3;\R)\), whose local gauge
transformations are STA rotors. The CSTA formulation extends the
internal gauge algebra,
\[
    \Cl^2(1,3;\R)\longrightarrow\Cl^2(2,4;\R),
\]
from six to fifteen generators, adding translations, dilations and
special conformal transformations.

\subsection{The Lorentz Sector}

The natural restriction
\[
    \Om(a,x)=\frac12\omega^{ab}(a,x)\,\e_a\wedgeprod\e_b,
    \qquad
    e^a=f^a=b=0,
\]
embeds a Lorentz-valued connection in CSTA. Its curvature stays in the
Lorentz sector, and Lorentz rotors preserve the restriction. A general
\(\Spin(2,4)\) transformation, however, carries a Lorentz connection
out of the sector, into all of \(P\), \(K\) and \(D\). The restriction
is therefore a gauge choice, not a gauge-invariant condition.

The Lorentz sector of CSTA is related to STA's bivectors (in the
opposite overall signature) by \(M_{ab}\mapsto-M_{ab}\), a Lie algebra
homomorphism: the structure constants differ by exactly \(-1\).

\subsection{The Translation Sector as Coframe}
\label{sec:coframe}

The translation sector is the soldering form. Identify it with an
orthonormal coframe,
\[
    e(\partial_\nu)=e^a{}_\nu\,\e_a,
    \qquad
    g_{\mu\nu}=e(\partial_\mu)\cdot e(\partial_\nu),
    \qquad
    E=e^{-1},
\]
and put the coframe's Levi-Civita spin connection in the Lorentz
sector. From the Christoffel symbols
\(\Gamma_{\lambda\mu\nu}=\tfrac12(\partial_\mu g_{\lambda\nu}
+\partial_\nu g_{\lambda\mu}-\partial_\lambda g_{\mu\nu})\), the frame
connection with lowered first index is
\[
    W(c,\mu,b)=\Gamma(E_c,\partial_\mu,E_b)
    +e_c\cdot e(\partial_\mu E_b),
\]
and the bivector connection is
\[
    \Om_M(\partial_\mu)=\tfrac12\,W(\e_k,\mu,\e_l)\,\e^k\wedgeprod\e^l .
\]
The normalization is forced: since \(\tfrac12[B,v]=B\cdot v\),
\(\Om_M(\partial_\mu)\cross v\) is the frame covariant derivative of
\(v\). Then:

\begin{itemize}
    \item The curvature of \(\Om_M\) is the Riemann bivector map,
    \(\F^{ab}{}_{\mu\nu}=R^{ab}{}_{\mu\nu}\). For Schwarzschild in
    Painlev\'e--Gullstrand form, Ricci vanishes and the Kretschmann
    scalar is \(48M^2/r^6\).

    \item With the full connection
    \(\Om=\Om_M+e\wedgeprod\ninf\), the translation sector of the
    curvature is the torsion, and it vanishes:
    \(\Pi_P\F=0\).
\end{itemize}

Where a GTG field is the identity plus a dyad \(N\) with \(N^2=0\),
the frame is immediate: \(e=\mathrm{Id}+N\Rightarrow E=\mathrm{Id}-N\).
This covers the Painlev\'e--Gullstrand and Kerr--Schild forms below.

\subsection{Example Geometries}

The following coframes serve as examples throughout.

\begin{center}
\renewcommand{\arraystretch}{1.2}
\begin{tabular}{>{\raggedright\arraybackslash}p{0.2\textwidth}
                >{\raggedright\arraybackslash}p{0.44\textwidth}
                >{\raggedright\arraybackslash}p{0.26\textwidth}}
    \toprule
    Geometry & Coframe & Character\\
    \midrule
    Schwarzschild & \(e^0=dt,\ e^i=dx^i+\beta^i dt,\
        \beta=\sqrt{2M/r}\,\hat x\) & vacuum; regular at \(r=2M\)\\
    Growing mass & as above, \(M=M(t)\) increasing & torsion-free,
        with matter\\
    Friedmann--Robertson--Walker (FRW) & \(e^0=dt,\ e^i=a(t)\,dx^i,\
        a=t^{2/3}\) & conformally flat\\
    de Sitter & \(e^0=dt,\ e^i=e^{Ht}dx^i\) & Einstein, \(\Lambda=3H^2\);
        conformally flat\\
    Static lapse & \(e^0=N(\mathbf x)\,dt,\ e^i=dx^i\) & Weyl, Ricci,
        Cotton nonzero\\
    Kerr & \(e=\mathrm{Id}-\phi\,l\,(l\cdot\ )\) (below) & vacuum,
        rotating\\
    \bottomrule
\end{tabular}
\end{center}

\subsection{Kerr in the CSTA}

GTG's position gauge field for Kerr, in Cartesian Kerr--Schild form,
is the identity plus a null dyad,
\[
    e=\mathrm{Id}-\phi\,l\,(l\cdot\ ),
    \qquad
    \phi=\frac{Mr^3}{r^4+a^2z^2},
    \qquad
    E=\mathrm{Id}+\phi\,l\,(l\cdot\ ),
\]
with \(l=\partial_t-\hat n\) ingoing and null in the flat metric, so
that \(g=\eta-2\phi\,l\,l\). As a conformal connection (coframe in the
translation sector, Levi-Civita connection in the Lorentz sector),
Kerr is torsion-free, Ricci-flat, has vanishing special-conformal and
dilation curvature, and its normal special-conformal field
(Section~\ref{sec:normal}) is \(f=0\). Its curvature density
\(2\grade{\F\rev\F}{0}\) is the Kretschmann scalar
\[
    48M^2\,\frac{r^6-15r^4u+15r^2u^2-u^3}{(r^2+u)^6},
    \qquad
    u=\frac{a^2z^2}{r^2}.
\]
GTG's position gauge field and its Kerr--Schild rotation gauge field
(in closed form) coincide exactly with the coframe and the Lorentz
connection held by the CSTA's translation and Lorentz sectors.

\subsection{GTG Gauge Freedom as CSTA Gauge Transformations}

GTG has two gauge freedoms: displacements of the position gauge field
and rotations. For Schwarzschild, the Newtonian
(Painlev\'e--Gullstrand) gauge,
\[
    e^0=dt,\qquad e^i=dx^i+\beta^i\,dt,\qquad
    \beta=\sqrt{2M/r}\,\hat x,
\]
and the Kerr--Schild gauge (Kerr at \(a=0\)) are related by the
displacement
\[
    t_{\mathrm{KS}}=t_{\mathrm{PG}}+F(r),\qquad
    F(r)=4M\ln\frac{\sqrt r+\sqrt{2M}}{\sqrt{2M}}-2\sqrt{2Mr},
\]
followed by a radial boost of rapidity \(\ln\bigl(1+\sqrt{2M/r}\bigr)\).
Pulled back by the displacement, the Kerr--Schild coframe is the
boosted Painlev\'e--Gullstrand coframe, and the pulled-back full CSTA
connection (Lorentz and translation sectors) is the
Painlev\'e--Gullstrand connection transformed by the boost rotor \(R\)
with the law of Section~5,
\[
    \Om'=R\,\Om\,\rev R-2(\nabla R)\rev R,
\]
outside and inside the horizon.

In the Painlev\'e--Gullstrand frame, the rotation gauge field is Doran
and Lasenby's closed form: \(\Om(E_0)=0\), and \(\Om(E_k)\) is a pure
boost taking \(\e_0\) to \(-B_{ik}\,\e_i\), with
\[
    B_{ik}=\sqrt{\frac{2M}{r^3}}\Bigl(\delta_{ik}
    -\tfrac32\,\hat r_i\hat r_k\Bigr).
\]

% --------------------------------------------------
\section{The Normal Conformal Connection and Weyl Rescaling}
\label{sec:normal}
% --------------------------------------------------

\subsection{The Normal Conditions}

Given the coframe in the translation sector and its Levi-Civita
connection in the Lorentz sector, the special-conformal field \(f\) is
fixed by the normal conditions
\[
    \Pi_P\F=0,\qquad
    \Pi_D\F=0,\qquad
    \Pi_M\F\ \text{trace-free},
\]
where the trace is the Ricci map
\(\Ric(\nu)=\sum_a\e^a\cdot\Pi_M\F(E_a,\partial_\nu)\). The conditions
are affine in the sixteen components of \(f\) and involve no
derivative of \(f\), so they hold pointwise. The system has full
rank \(16\), and its unique solution is the Schouten tensor,
\[
    \boxed{
        f=P=\frac12\left(\Ric-\frac{R}{6}\,e\right),
    }
\]
with coefficient exactly \(1\) in this normalization. For Ricci-flat
spacetimes \(f=0\). This is the normal conformal Cartan
connection,
\[
    \Om_{\mathrm{normal}}=\Om_M+e\wedgeprod\ninf+P\wedgeprod\nzero .
\]

\subsection{Curvature of the Normal Connection}

With \(f=P\):
\[
    \Pi_P\F=0,\qquad
    \Pi_D\F=0,\qquad
    \Pi_M\F=C\ \text{(Weyl)},\qquad
    \Pi_K\F=\D P\ \text{(Cotton)}.
\]
The Cotton sector vanishes for FRW (conformally flat) and in vacuum,
and is nonzero for the static-lapse metric. The normal connection is
flat exactly when spacetime is conformally flat.

\subsection{The Quadratic Curvature Density}

The scalar contraction
\[
    \mathcal{I}=g^{\mu\rho}g^{\nu\sigma}\,
    \grade{\F_{\mu\nu}\rev{\F}_{\rho\sigma}}{0}
\]
of the full fifteen-component curvature satisfies
\[
    \boxed{
        \mathcal{I}\bigl[\Om_{\mathrm{normal}}\bigr]
        =\frac12\,C_{abcd}C^{abcd},
    }
\]
the Weyl-squared density, whereas for the Lorentz connection alone it
is \(\tfrac12R_{abcd}R^{abcd}\) in vacuum and also sees the Ricci
tensor otherwise. A quadratic curvature action of the normal
conformal connection is therefore Weyl-squared gravity. Its field equations are in
Section~\ref{sec:dynamics}.

\subsection{Weyl Rescaling as a Gauge Transformation}

Under a Weyl rescaling of the coframe, \(e\mapsto e^\lambda e\), with
\(\lambda\) an arbitrary function of all four coordinates, the normal
connection of the rescaled coframe is the normal connection of \(e\)
transformed by the versor
\[
    \boxed{
        V=\exp\left(-\tfrac{\lambda}{2}\,\nzero\wedgeprod\ninf\right)
          \exp\left(\tfrac12\,\Upsilon\wedgeprod\nzero\right),
        \qquad
        \Upsilon=\nabla\lambda\ \text{(in the frame of }e\text{)},
    }
\]
with the transformation law of Section~5. The dilation factor alone
leaves the Lorentz sector unchanged, scales the translation sector by
\(e^{\lambda}\) and the special-conformal sector by \(e^{-\lambda}\),
and switches on the dilation sector with
\[
    b=d\lambda,
\]
while the curvature's dilation sector stays zero; the special
conformal factor removes \(b\) again. The Weyl-squared density scales
as \(C^2\mapsto e^{-4\lambda}C^2\).

In this sense the dilation and special-conformal sectors carry a
choice of scale, a gauge, rather than additional degrees of freedom:
a pure-gauge dilation field \(b=d\lambda\) is curl-free, and the
special-conformal field is fixed algebraically by the coframe.

% --------------------------------------------------
\section{Gravitational Dynamics}
\label{sec:dynamics}
% --------------------------------------------------

The connection and curvature establish the kinematic structure of the
theory; a gravitational theory additionally requires an action. The
quadratic curvature action of the normal connection is Weyl-squared
gravity, and recovering Einstein gravity requires a compensator that
breaks the conformal symmetry.

\subsection{Field Equations}

For a density \(\mathcal L(e,\partial e,\partial\partial e)\) of
order at most two in the coframe, the Euler--Lagrange expression
\[
    E[a,\mu]=\frac{\delta S}{\delta e^a{}_\mu}
    =A-\partial_\mu A^\mu+\sum_{\mu\le\nu}\partial_\mu\partial_\nu
    A^{\mu\nu},
\]
with \(A,A^\mu,A^{\mu\nu}\) the coefficients of \(h\), \(\partial_\mu
h\) and \(\partial_\mu\partial_\nu h\) in the first variation, is
local and can be evaluated pointwise. Densities carry the volume
\(\sqrt{g}=|\det e|\), read from the outermorphism of the coframe on
the pseudoscalar of directions. The two basic densities give
\begin{align*}
    \sqrt g\,R:&\quad
    E[a,\mu]=-2\sqrt g\,e_{a\nu}G^{\mu\nu},
    &&\text{i.e.}\quad
    \frac{\delta S}{\delta g_{\mu\nu}}=-\sqrt g\,G^{\mu\nu},\\
    \sqrt g\,C^2:&\quad
    E[a,\mu]=-8\sqrt g\,e_{a\nu}B^{\mu\nu},
    &&\text{i.e.}\quad
    \frac{\delta S}{\delta g_{\mu\nu}}=-4\sqrt g\,B^{\mu\nu},
\end{align*}
with \(G\) the Einstein tensor and
\[
    B_{ab}=\nabla^c\nabla_cP_{ab}-\nabla^c\nabla_aP_{bc}
    +P^{cd}C_{acbd}
\]
the Bach tensor. Schwarzschild,
Kerr, FRW and de Sitter are Bach-flat. Both densities satisfy the
Noether identity of local Lorentz invariance
(\(\eta E e^{\mathsf T}\) symmetric), and for Weyl squared the Weyl
trace \(E[a,\mu]\,e^a{}_\mu\) vanishes, since \(B\) is trace-free.

\subsection{The Scale Tractor as Compensator}

A choice of scale \(\sigma>0\), picking out the metric \(g/\sigma^2\),
defines a field of \(\R^{2,4}\) vectors,
\[
    \boxed{
        \mathcal{I}_\sigma=\sigma\,\ninf+\nabla\sigma+\rho\,\nzero,
        \qquad
        \rho=-\tfrac14\left(\Box\sigma+J\sigma\right),
        \qquad
        J=\tfrac{R}{6}=\operatorname{tr}P,
    }
\]
the scale tractor, moved by the normal connection as a
conjugation-type field,
\[
    \D_a\mathcal I_\sigma=a\cdot\nabla\mathcal I_\sigma
    +\Om_{\mathrm{normal}}(a)\cross\mathcal I_\sigma .
\]
It is the CSTA form of Penrose's infinity twistor. Its properties:

\begin{itemize}
    \item \textbf{Weyl covariance.} Under \(e\mapsto e^\lambda e\),
    \(\sigma\mapsto e^\lambda\sigma\), it transforms by the versor of
    the Weyl rescaling,
    \(\mathcal I\mapsto V\mathcal I\rev V\), so
    \[
        \mathcal I_\sigma\cdot\mathcal I_\sigma
        =|\nabla\sigma|^2+2\sigma\rho
    \]
    is Weyl-invariant.

    \item \textbf{Einstein--Hilbert density.} At \(\sigma=1\),
    \[
        \mathcal I\cdot\mathcal I=-\frac J2=-\frac{R}{12},
    \]
    so the conformally invariant density
    \(\sqrt g\,\sigma^{-4}\,\mathcal I_\sigma\cdot\mathcal I_\sigma\)
    is \(-\sqrt g\,R/12\) in the gauge \(\sigma=1\).

    \item \textbf{Parallel exactly on Einstein metrics.}
    \(\D\mathcal I_\sigma=0\) exactly when \(g/\sigma^2\) is Einstein.
    For Ricci-flat metrics \(\mathcal I=\ninf\) (null) at \(\sigma=1\);
    for de Sitter with Hubble rate \(H\),
    \(\mathcal I\cdot\mathcal I=H^2=\Lambda/3\). A non-null
    compensator is a cosmological constant.

    \item \textbf{Integrability.} A parallel \(\mathcal I\) satisfies
    \(\F(a,b)\cross\mathcal I=0\). The null space of these conditions
    always contains \(\nzero\) when the torsion sector vanishes, but
    \(\nzero\) has scale \(\mathcal I\cdot\nzero=0\); where it is the
    only solution (as for the static lapse and the growing mass), no
    Einstein scale exists. For Schwarzschild the null space is
    \(\mathrm{span}(\nzero,\ninf)\).
\end{itemize}

\subsection{Second-Order Form: The Compensator Action}

The Weyl-invariant action
\[
    \boxed{
        S_c=\int\sqrt{g}\,\sigma^{-4}
        \bigl(\mathcal I_\sigma\cdot\mathcal I_\sigma+c\bigr)\,d^4x,
    }
\]
equal to \(-\tfrac1{12}\int\sqrt{g'}\,(R'+12c)\) with
\(g'=g/\sigma^2\), is varied in the coframe and in \(\sigma\).

\begin{itemize}
    \item \textbf{Einstein's equations.} In the gauge \(\sigma=1\), the
    coframe equation is
    \(\sqrt g\,e_{a\nu}\bigl(\tfrac16G^{\mu\nu}+c\,g^{\mu\nu}\bigr)=0\),
    that is,
    \[
        \boxed{
            G_{\mu\nu}+6c\,g_{\mu\nu}=0,
            \qquad
            \Lambda=-6c,
        }
    \]
    and on shell \(\mathcal I\cdot\mathcal I=-R/12=-2c\).

    \item \textbf{The scale equation is the trace.} Off shell,
    \(E_\sigma\,\sigma+E[a,\mu]\,e^a{}_\mu=0\): the Noether identity of
    Weyl invariance. The scale carries no equation of its own, and the
    gauge \(\sigma=1\) loses nothing.

    \item \textbf{Solutions.} de Sitter with \(c=-H^2/2\), Kerr with
    \(c=0\), and a Weyl-rescaled Schwarzschild together with its scale
    \(\sigma=e^\lambda\) solve both equations; de Sitter with \(c=0\),
    and the rescaled Schwarzschild with \(\sigma=1\), do not. Solutions
    are classes under Weyl rescaling together with their scale.
\end{itemize}

With the infinity twistor as compensator, this action is a
reformulation of general relativity with a cosmological constant, not
a new theory.

\subsection{First-Order Form: The Full Connection Varied}

Two constant compensators, the origin \(X=\nzero\) and the infinity
twistor \(\mathcal I=\ninf\) (null, \(X\cdot\mathcal I=1\)), break
\(\Spin(2,4)\) to the Lorentz group. Their covariant derivatives are
the coframe and the special-conformal field,
\[
    \D X=\Om\cross\nzero,
    \qquad
    \D\mathcal I=\Om\cross\ninf,
\]
which reproduce the translation and special-conformal sectors. Writing
\(\langle\cdots\rangle\) for the pseudoscalar coefficient of a
pseudoscalar-valued 4-form evaluated on the volume of directions,
define
\begin{align*}
    S_1&=\bigl\langle \F\wedge\D X\wedge\D X\wedge X\wedge\mathcal I
        \bigr\rangle,\\
    S_2&=\bigl\langle \D\mathcal I\wedge\D\mathcal I\wedge\D X\wedge\D X
        \wedge X\wedge\mathcal I\bigr\rangle,\\
    S_3&=\bigl\langle \D X\wedge\D X\wedge\D X\wedge\D X
        \wedge X\wedge\mathcal I\bigr\rangle,
\end{align*}
a conformal Stelle--West form: \(S_1\) is Einstein--Cartan plus a term
linear in the special-conformal field, \(S_2\) makes that field
algebraic, and \(S_3\) is a cosmological term. The action
\[
    \boxed{
        S=S_1+\beta S_2+\gamma S_3
    }
\]
is varied in all \(60\) components of \(\Om\), with no other field.
Its solutions have
\begin{itemize}
    \item a torsion-free Lorentz sector (a constant contorsion added
    to it is not a solution);
    \item an algebraic special-conformal field, \(f=e/\beta\);
    \item an Einstein metric with
    \[
        \boxed{
            \Lambda=6\left(\gamma-\frac1\beta\right);
        }
    \]
    \item a dilation field that drops out of the action entirely
    (Schwarzschild remains a solution with a \(b\) that has curl).
\end{itemize}
Along the family (Levi-Civita, \(e\), \(f=ke\), \(b=0\)) each term's
equations are polynomial in \(k\), and for an Einstein metric they
vanish exactly when \(k=1/\beta\) and \(\gamma=k+\Lambda/6\). Thus
Schwarzschild and Kerr solve at \(\beta=4\), \(\gamma=\tfrac14\), and
de Sitter at \(\beta=2\), \(\gamma=\tfrac12+\tfrac{H^2}2\). The static
lapse, which is not Einstein, has no solution for any
\(k,\beta,\gamma\).

\textbf{Normal connection.} With \(\gamma=0\), the solutions are
exactly the normal conformal connections of Einstein metrics: de
Sitter's Schouten tensor is \(-(\Lambda/6)\,e\) in this sign of \(R\),
and \(\Om_{\mathrm{normal}}\) solves at \(\beta=-6/\Lambda\). The
cosmological constant then comes from the coupling of the
special-conformal field, and the Ricci-flat normal connection is the
limit \(\beta\to\infty\).

Einstein's equations with a cosmological constant thus follow from
varying the full CSTA connection alone, once the origin and the
infinity twistor are fixed.

\subsection{The Action Quadratic in Curvature: Weyl Gravity}

The MacDowell--Mansouri form with the same compensators,
\[
    \boxed{
        S_{\mathrm{MM}}=\bigl\langle\F\wedge\F\wedge X\wedge\mathcal I
        \bigr\rangle,
    }
\]
retains only the Lorentz sector of \(\F\), giving Kaku, Townsend and
van Nieuwenhuizen's \(\epsilon_{abcd}\F^{ab}\wedge\F^{cd}\). Because
\([P,P]=0\) in the conformal algebra, \(\Pi_M\F\) contains an
\(e\wedge f\) term but no \(e\wedge e\) term, and, unlike de Sitter
MacDowell--Mansouri gravity, no Einstein--Hilbert term arises.

\begin{itemize}
    \item \textbf{Density.} At the normal connection,
    \(S_{\mathrm{MM}}=-\sqrt g\,C^2\): minus the Weyl-squared action.

    \item \textbf{Special-conformal equation.} It holds at the normal
    connection identically, so it forces \(f=P\).

    \item \textbf{Lorentz equation.}
    \(4\,\epsilon_{abcd}\,e^c\wedge(\Pi_K\F)^d=0\), with
    \(\Pi_K\F=\D P\) the Cotton 2-form: the metric must be Cotton-free
    (divergence-free Weyl tensor) in the scale fixed by the
    compensators.

    \item \textbf{Dilation equation.} Identically zero; \(b\) drops
    out.

    \item \textbf{Solutions.} The normal connections of Schwarzschild,
    Kerr, de Sitter and FRW solve all equations. FRW is not Einstein,
    so the theory is not general relativity. FRW with \(f=0\) does not
    solve it, nor does a Weyl-rescaled Schwarzschild (Bach-flat, but
    with nonzero Cotton tensor): the fixed compensators break the Weyl
    invariance of Weyl gravity.
\end{itemize}

The symmetry breaking by this split therefore gives Weyl gravity with
a Cotton condition, not Einstein gravity. The term that makes the
difference is \(S_2\), quadratic in the special-conformal field.

\subsection{Quadratic Admixture: Einstein Solutions and a Ghost}

\textbf{First order.} For \(S_1+\beta S_2+\gamma S_3+\mu S_{\mathrm{MM}}\),
Einstein metrics still solve the equations, with the
special-conformal field shifted,
\[
    f=k\,e,\qquad
    k=\frac{1-4\mu\gamma}{\beta},\qquad
    \Lambda=6(\gamma-k).
\]

\textbf{Second order.} With \(\alpha\sqrt g\,C^2\) added to the
compensator action, Einstein metrics remain solutions (they are
Bach-flat), and so does a Weyl-rescaled Schwarzschild with its scale,
since \(\sqrt g\,C^2\) is Weyl-invariant.

\textbf{Linearized modes.} For a transverse-traceless plane wave
\(e=\mathrm{Id}+\varepsilon H\cos(k\cdot x)\) about flat space, with
\(k^2=\omega^2-q^2\),
\[
    \sqrt g\,R:\ -2k^2H,
    \qquad
    \sqrt g\,C^2:\ 4k^4H,
\]
and both vanish for null waves. The compensator action (whose
Einstein--Hilbert coefficient is \(-\tfrac1{12}\)) plus
\(\alpha\sqrt g\,C^2\) therefore has
\[
    \boxed{
        L=\frac{k^2}{6}\bigl(1+24\alpha k^2\bigr),
        \qquad
        \text{propagator}\ \propto\
        \frac1{k^2}-\frac1{k^2-m^2},
        \qquad
        m^2=-\frac1{24\alpha}:
    }
\]
a massive spin-2 mode with residue opposite to the graviton's,
Stelle's ghost, non-tachyonic only for \(\alpha<0\). The scale adds no
mode, since its equation is the trace of the coframe equation, so the
compensator does not remove the ghost. With \(\alpha=0\) the theory is
general relativity; any Weyl-squared admixture keeps every Einstein
solution and adds the massive ghost.

% --------------------------------------------------
\section{Twistors}
\label{sec:twistors}
% --------------------------------------------------

\subsection{The Complex Structure of the Even Subalgebra}

Because the pseudoscalar \(I\) is central in the even subalgebra and
\(I^2=-1\), it acts there as the imaginary unit, and
\[
    \boxed{
        \G^{+}\cong M(4,\C),
        \qquad
        \Spin(2,4)\cong SU(2,2).
    }
\]
The complex numbers of twistor theory are not added to the real
algebra; they are its pseudoscalar.

\subsection{Twistor Space as a Minimal Left Ideal}

The element
\[
    P_{\T}=\frac14\left(1+I\,\e_x\e_y\right)\left(1+I\,\e_t\e_+\right)
\]
is a primitive idempotent and its own reverse,
\(P_{\T}^2=P_{\T}=\rev{P}_{\T}\). Twistor space is the minimal left
ideal
\[
    \boxed{
        \T=\G^{+}P_{\T},
    }
\]
of four complex (eight real) dimensions. Flipping the signs in
\(P_{\T}\) gives four orthogonal idempotents summing to one, so every
even field \(\Psi\) is four twistors side by side. The transformation
law \(\Psi'=V\Psi\) acts on \(\T\) unchanged, \(Z'=VZ\), and so do the
covariant derivative, transport and holonomy.

The twistor norm
\[
    h(Z_1,Z_2)=4\,\grade{\rev{Z}_1Z_2}{0,6},
\]
read as a complex number (scalar plus \(I\) part), is Hermitian with
signature \((2,2)\). Reversion is its adjoint, so versors preserve it:
they act as matrices of unit determinant in \(SU(2,2)\), and bivectors
as traceless \(h\)-antihermitian matrices. The complex trace on \(\T\)
is \(4\grade{A}{0,6}\).

\subsection{The \texorpdfstring{\(\omega/\pi\)}{omega/pi} Split}

The null basis splits a twistor into two halves,
\[
    \omega=\tfrac12\,\ninf\nzero Z,
    \qquad
    \pi=\tfrac12\,\nzero\ninf Z,
    \qquad
    \ninf\omega=0,\quad\nzero\pi=0,
\]
by complementary idempotents. Each half is \(h\)-null and the two pair
with each other; in an adapted basis
\(h=2\,\mathrm{Re}(\bar\omega\cdot\pi)\). The generators act on the
split as follows:

\begin{center}
\begin{tabular}{ll}
    \toprule
    Generator & Action on \((\omega,\pi)\)\\
    \midrule
    Lorentz, \(\e_a\wedgeprod\e_b\) & turns each half separately\\
    Translation, \(\e_a\wedgeprod\ninf\) & maps \(\pi\) into \(\omega\)\\
    Special conformal, \(\e_a\wedgeprod\nzero\) & maps \(\omega\) into
        \(\pi\)\\
    Dilation, \(\nzero\wedgeprod\ninf\) & scales the halves oppositely\\
    \bottomrule
\end{tabular}
\end{center}

\subsection{Incidence and Light Rays}

A twistor and an event are incident when the event's conformal point
annihilates the twistor:
\[
    \boxed{
        XZ=0
        \quad\Longleftrightarrow\quad
        \nzero\rev{T}_xZ=0
        \quad\Longleftrightarrow\quad
        Z=T_x\pi_0 ,
    }
\]
with \(\pi_0\) in the \(\pi\) half: translated back to the origin, the
twistor passes through the origin. Equivalently
\[
    \omega=\tfrac12\,x\,\ninf\,\pi,
\]
the index-free form of Penrose's \(\omega=ix\pi\); in spinor
coordinates \(\omega=-\tfrac{i}{\sqrt2}\,x^{AA'}\pi\), with \(x^{AA'}\)
Hermitian of determinant \(x^2\). The incident twistors of an event
form a complex 2-plane (its \(\alpha\)-plane), and the
\(\alpha\)-planes of two events meet in a complex line exactly when
the events are null separated: the Klein correspondence.

For a null twistor, \(h(Z,Z)=0\), the real vectors annihilating \(Z\)
form a totally null 2-plane, the light ray of \(Z\). A non-null
twistor is incident with no real event.

\subsection{Momentum and Null Flag}

The momentum of a twistor is the moment map of the translations,
\[
    p^a=-\mathrm{Im}\,h\bigl(Z,(\e^a\wedgeprod\ninf)Z\bigr)
    =4\,\Bigl\langle(\e^a\wedgeprod\ninf)\,
        \grade{Z\rev Z}{4}\,I\Bigr\rangle_0 ,
\]
the second form following from the cyclicity of the scalar part and
the fact that \(Z\rev Z\), being its own reverse, has only grades 0
and 4. It is null and future-pointing, Penrose's \(p=\pi\bar\pi\),
independent of \(\omega\) and of the phase of \(Z\), and quadratic in
\(Z\). The helicity is \(s=\tfrac12h(Z,Z)\). Conversely, a twistor
through any event with any given null momentum can be constructed.

The phase of \(Z\) is seen by the null flag, the real part of the
self-dual bivector \(\epsilon_{AB}\pi^A\pi^B\). No even element
conjugates a twistor, since \(I\) is central in \(\G^+\), but the odd
element \(C=\e_t\e_y\e_z\) anticommutes with \(I\) and supplies
\(\epsilon_{AB}\); the flag is the Lorentz part of
\[
    \sqrt8\,\grade{\ninf\,\pi\,C\,\rev\pi}{2},
\]
a null bivector \(k\wedgeprod s\) with \(s\) a unit spacelike
direction across the ray, defined up to multiples of \(k\).
Multiplying \(Z\) by \(e^{I\theta}\) turns \(s\) by \(2\theta\) about
\(k\), and a Lorentz rotor carries the flag of \(Z\) to that of \(RZ\)
as it carries vectors.

\subsection{Local Twistor Transport}

The covariant derivative acting on twistors,
\[
    \D_aZ=a\cdot\nabla Z+\tfrac12\Om(a,x)Z,
\]
is Penrose's local twistor transport when \(\Om\) is the normal
conformal connection. Its curvature,
\([\D_a,\D_b]Z=\tfrac12\F(a,b)Z\), carries the Weyl tensor in the
Lorentz sector and the Cotton tensor in the special-conformal sector,
and vanishes exactly when spacetime is conformally flat: it vanishes
for FRW and not for the static lapse. Transport maps twistors
by elements of \(SU(2,2)\), and small holonomy loops equal
\(\exp(-\tfrac12\epsilon^2\F)\) to \(O(\epsilon^3)\).

\textbf{Null geodesics.} Along a null geodesic, the event \(x\) and
the local twistor \(Z\) of its light ray evolve by
\[
    \dot x=E(k),
    \qquad
    \dot Z=-\tfrac12\,\Om_{\mathrm{normal}}(\dot x)\,Z,
    \qquad
    k=p(Z),
\]
where the translation part of the connection along \(\dot x\) is
\(k\wedgeprod\ninf\). The Lorentz part turns \(\pi\) as parallel
transport, so \(k\) obeys the geodesic equation in the frame and stays
exactly null; the translation part moves \(\pi\) into \(\omega\) at
the rate the event moves, so a twistor through the current event
stays through it (\(\omega=0\)). The geodesic equation is thus local
twistor transport by the CSTA's own connection, and the null flag
carries the transverse screen along the ray. Equivalently, the
transport can be accumulated as a versor, \(Z=UZ_0\) with
\(\dot U=-\tfrac12\Om_{\mathrm{normal}}(\dot x)\,U\) and
\(U\in\Spin(2,4)\).

\subsection{The Complex Weyl Invariant}

Lorentz bivectors preserve the \(\omega/\pi\) split and act on the
\(\omega\) half by their self-dual part. The trace of the curvature
squared over the \(\omega\) half is therefore complex:
\[
\boxed{
    \W=
    2\left\langle\tfrac12\,\ninf\nzero\,
    \F(\e_a,\e_b)\,\rev{\F}(\e^a,\e^b)
    \right\rangle_{0,6}
    =
    \frac12\left(C_{abcd}C^{abcd}+i\,C_{abcd}{}^{*}C^{abcd}\right),
}
\]
with \(I\) read as \(i\). For Kerr,
\[
    \W=\frac{24M^2}{(r-ia\cos\theta)^6},
\]
and for Schwarzschild it is real. Its real part is the quadratic
density of Section~\ref{sec:normal}; its imaginary part is half the
Chern--Pontryagin density, a total derivative. A complex action
\(\int\W\) therefore has the Weyl-squared action as its real part and
a topological term as its imaginary part, which leaves the field
equations unchanged.

The twistor holonomy of Kerr shows the same structure: in Petrov type
D the Weyl curvature acts on self-dual bivectors as multiplication by
\(\Psi_2\propto(r-ia\cos\theta)^{-3}\), and the boost (gravitomagnetic)
fraction of small-loop holonomies follows
\(|\arg\Psi_2|=3\arctan(a\cos\theta/r)\): zero on the equator and for
\(a=0\), largest near the poles.

% --------------------------------------------------
\section{Dirac Matter}
\label{sec:dirac}
% --------------------------------------------------

Gauge Theory Gravity is distinguished from general relativity by its
coupling to Dirac fields: the spin of the matter sources torsion
\cite[Sections 3--4]{LDG}. This section carries Lasenby, Doran and Gull's Dirac
action into the first-order form of Section~\ref{sec:dynamics} and
checks it against an independent implementation of their formulas.

\subsection{Spinors as Twistors}

A Dirac spinor is a twistor \(Z\in\G^+P\). The Dirac algebra acts
through
\[
    \boxed{
        \Gamma_a=I\,\e_a\,p,
        \qquad
        p=\tfrac1{\sqrt2}\left(\nzero+\ninf\right)
        =\tfrac1{\sqrt2}\left(X+\mathcal I\right),
    }
\]
which satisfy \(\{\Gamma_a,\Gamma_b\}=2\eta_{ab}\) and
\(\Gamma_a\Gamma_b=\e_a\e_b\) for \(a\ne b\). The Lorentz sector
therefore acts on \(Z\) as GTG's rotation gauge field acts on its
spinor, and the twistor norm is the Dirac bar pairing,
\(h(Z_1,Z_2)=\bar\psi_1\psi_2\). The other generators act
algebraically:
\[
    \nzero\wedge\ninf=\gamma_5,
    \qquad
    P_a=-\sqrt2\,I\,\Gamma_a\Pi_+,
    \qquad
    K_a=-\sqrt2\,I\,\Gamma_a\Pi_-,
    \qquad
    \Pi_\pm=\tfrac12(1\pm\gamma_5),
\]
so the \(\omega\) and \(\pi\) halves of a twistor are its Weyl
spinors. The map \(Z=\psi\,Z_0\) from the even spacetime algebra, with
\(Z_0\) the unit twistor fixed by \(\Gamma_0Z_0=Z_0\) and
\(\e_y\e_xZ_0=IZ_0\), carries \(\gamma_\mu\psi\gamma_0\mapsto\Gamma_\mu Z\),
\(\psi I\sigma_3\mapsto IZ\), and every bilinear. In particular
\(\grade{\psi\rev\psi}{0}\mapsto h(Z,Z)\) and
\(\grade{\gamma^a\phi I\gamma_3\rev\psi}{0}\mapsto
\operatorname{Re}h(Z_\psi,I\Gamma^aZ_\phi)\).

\subsection{The Action}

The Dirac density of a connection is
\[
    \boxed{
        \mathcal L_D=|\det e|\left(
        \operatorname{Re}h\left(Z,I\Gamma^a\D_{E_a}Z\right)-m\,h(Z,Z)
        \right),
    }
\]
with the coframe read from \(\D X\) and \(E_a\) its frame. The
transport \(\D\) is either the full twistor transport (all fifteen
generators acting on \(Z\)) or its Lorentz sector, the stabilizer of
\(X\) and \(\mathcal I\). The Lorentz coupling is Lasenby, Doran and
Gull's \cite[(3.56)]{LDG}. The full transport adds the \(h\)-Hermitian
part of \(\tfrac12I\Gamma^a(\Om-\Om_M)(E_a)\):
\[
    \sqrt2\,\sigma
    +\tfrac{\sqrt2}4\left(\operatorname{tr}F\,\sigma+F^{ab}T_{ab}\right),
    \qquad
    \sigma=\bar\psi\psi,
    \quad
    T_{ab}=\bar\psi\sigma_{ab}(1-\gamma_5)\psi,
    \quad
    F=fe^{-1},
\]
from the translation and special-conformal sectors; the dilation
sector adds nothing. For \(f=ke\) this is a parity-even mass made of
the compensators, \(m\mapsto m-\sqrt2(1+k)\). With LDG's \(R\), which
has the sign used here, their action
\(\det(h)^{-1}(\tfrac12R-\kappa L_D)\) is
\[
    \boxed{
        S_1+\beta S_2+\gamma S_3+4\kappa\,\mathcal L_D
    }
\]
up to a factor \(-4\).

\subsection{Field Equations}

\textbf{Dirac equation.} Varying \(Z\) gives
\[
    I\Gamma^a\D_aZ+\tfrac12I\Gamma(\tau)Z=mZ,
    \qquad
    \tau_b=-T^a{}_{ab},
    \qquad
    T(a,b)=-\nzero\rfloor\F_P(a,b),
\]
with \(\F_P\) the translation-sector curvature of the Lorentz and
translation parts of \(\Om\). (The dilation part would add \(b\wedge e\),
which the Dirac field does not see.) With the Lorentz coupling this is
LDG's coupled Dirac equation \cite[(3.62)]{LDG} component by
component. It is minimally coupled exactly when the torsion trace
vanishes.

\textbf{Sources.} The Dirac part of \(\delta S/\delta\Om\) is, sector by
sector:
\begin{itemize}
    \item Lorentz: \(|\det e|\,S\cdot g^\mu\), with
    \(S=\tfrac12\psi I\gamma_3\rev\psi\) and \(g^\mu=\bar h(e^\mu)\). This
    is LDG's spin tensor \(S\cdot\bar h(a)\).
    \item Translation: \(-|\det e|\,E^\mu{}_b\,\mathcal T^b{}_a\), with
    \(\mathcal T(a)=\grade{a\cdot\D\psi\,I\gamma_3\rev\psi}{1}-a\,L\), LDG's
    stress-energy \cite[(4.20)]{LDG}.
    \item Special conformal: zero for the Lorentz coupling; for the full
    transport \(|\det e|\,E^\mu{}_bC^b{}_a\), with
    \(C^b{}_a=\operatorname{Re}h(Z,\tfrac12I\Gamma^bK_aZ)\).
    \item Dilation: zero.
\end{itemize}

\textbf{Spin--torsion.} Levi-Civita plus LDG's contorsion,
\(\Om(\partial_\mu)\mapsto\Om(\partial_\mu)+\tfrac12\kappa\,g_\mu\cdot S\),
solves the Lorentz equations under either transport. The torsion is
\[
    \boxed{
        T(a,b)=\kappa\,(a\wedge b)\cdot S,
    }
\]
the spin density, and it is traceless. So the minimally coupled action
gives the minimally coupled equation, LDG's consistency
loop \cite[(4.32)--(4.35)]{LDG}.

\textbf{The special-conformal field.} Under the full transport, with
\(S_2=-2\sqrt g\left((\operatorname{tr}F)^2-F^a{}_bF^b{}_a\right)\), the
special-conformal equation is solved by
\[
    F^b{}_a=\frac1\beta\left(1+\frac\kappa3\operatorname{tr}C\right)
    \delta^b_a-\frac\kappa\beta\,C^b{}_a,
    \qquad
    \operatorname{tr}C=\sqrt2\,\sigma:
\]
the Dirac field moves \(f\) away from \(e/\beta\).

\textbf{The quadratic action.} Adding
\(\alpha\grade{\F\wedge\F\wedge X\wedge\mathcal I}{}\), whose
\(R\wedge R\) part is the Euler density, adds
\(\alpha\,\epsilon(T\wedge f-e\wedge\D f)\) to the Lorentz equations.
With the Lorentz coupling \(f\propto e\), and the term only rescales
\(S_1\)'s: the loop holds for every \(\alpha\). With the full transport
the Dirac-shifted \(f\) varies in space, and the torsion acquires a trace
linear in \(\alpha\). Dirac matter then requires \(\alpha=0\).

\subsection{A Self-Gravitating Solution}

Challinor, Lasenby, Doran and Gull \cite{CLDG} found a homogeneous Dirac
cosmology in GTG, extended with a cosmological constant by Leong
\cite{Leong}. In units \(G=1\), \(\kappa=8\pi\):
\[
    \bar h(b)=b+H(b\cdot x)\gamma_0,\qquad
    \psi=\sqrt\rho\,e^{I\beta/2}e^{-I\sigma_3\chi}L,\qquad
    3H^2=\Lambda+12\pi^2\rho^2+8\pi m\rho\cos\beta,
\]
with \(\cos\beta=\pm1\), \(\dot\rho=-3\rho H\) and
\(\dot\chi=m\cos\beta+3\pi\rho\), in closed form. The
\(12\pi^2\rho^2\) term is spin--torsion. In the CSTA the coframe
\(e^0=dt\), \(e^i=dx^i-Hx^i\,dt\) is identity plus a dyad. With
Levi-Civita plus LDG's contorsion, \(f=e/\beta\),
\(\Lambda=6(\gamma-1/\beta)\) and the Lorentz coupling, all
60 connection equations and the spinor equation hold. (A repulsive
\(\Lambda\) enters LDG's density as \(+\Lambda\det(h)^{-1}\) with
this sign of \(R\).)

Under the full transport the special-conformal equation shifts \(f\).
Eliminating it from \(12\operatorname{tr}F-2\beta((\operatorname{tr}F)^2
-F{:}F)+4\kappa\,C\cdot F\) leaves \(24/\beta\), which sits in
\(\Lambda\); a mass \(4\sqrt2\kappa\sigma/\beta\), which joins the
compensator mass; and a contact term
\[
    \frac{\kappa}{2\beta}\left(\tfrac13(\operatorname{tr}C)^2-C{:}C\right)
    =-\frac{\kappa\sigma^2}{24\beta}
\]
on this state. The solution keeps its form, with
\[
    m\to m_0=m-\sqrt2\left(1+\tfrac1\beta\right),\qquad
    12\pi^2\rho^2\to12\pi^2a\rho^2,\quad a=1+\frac2{9\beta},
\]
and the closed forms carry over with \(\sqrt a\,\alpha\) in place of
\(\alpha=\sqrt{3\Lambda}/2\); again every equation holds. The
compensator mass is absorbed into the bare mass, so the fermion mass
does not fix \(\beta\); \(\beta\) fixes the strength \(a\) of the
effective spin--torsion term.

% --------------------------------------------------
\section{Open Questions}
\label{sec:questions}
% --------------------------------------------------

\subsection{Fundamental Versor Fields}

Pure-gauge connections \(\Om=-2(\nabla V)\rev V\) are flat with
trivial holonomy, so versor fields alone describe only flat
connections. Can the gravitational degrees of freedom nevertheless be
represented by normalized CSTA versor fields together with independent
connection data, with the bivector-valued connection emerging from
their derivatives?

\subsection{Displacement Gauge Field}

The translation sector is an orthonormal coframe of a torsion-free
metric. Does it represent Doran and Lasenby's \(h(a)\), with its
fiducial flat background, in general? Open cases include Doran's
spinning Painlev\'e--Gullstrand gauge of Kerr. GTG's field equations
with Dirac matter hold pointwise (Section~\ref{sec:dirac}), and a
homogeneous Dirac cosmology solves them; localized self-gravitating
solutions are not constructed.

\subsection{The Dilation Field}

A pure-gauge dilation field \(b=d\lambda\) is curl-free and leaves the
curvature's dilation sector zero. A dilation field with nonzero curl
changes that sector. What dynamics, if any, does it carry?

\subsection{Einstein Gravity and Beyond}

With the infinity twistor as compensator the theory is general
relativity with a cosmological constant, and a Weyl-squared admixture
adds a massive spin-2 ghost of mass \(m^2=-1/(24\alpha)\). Open:
\begin{itemize}
    \item whether \(\alpha=0\), or the quadratic term is an effective
    correction below \(m\), without Dirac matter or with the Lorentz
    coupling (with the full transport, Dirac matter requires
    \(\alpha=0\));
    \item the scalar and vector sectors of the linearization, and the
    linearization of the first-order action;
    \item whether every Cotton-free solution of the coframe equation of
    \(S_{\mathrm{MM}}\) is Bach-flat;
    \item dynamics for the compensators, which are fixed in both the
    second- and first-order forms;
    \item the physical meaning of \(\beta\);
    \item matter other than Dirac's, the general form of the contact
    term left by the Dirac-shifted special-conformal field (checked
    only on the homogeneous cosmology), and the conformal weight of
    the Dirac field.
\end{itemize}

\subsection{Twistors}

\textbf{Twistor form of the dynamics.} Can the field equations be
expressed on \(\T\)? A global twistor space exists for conformally flat
spacetimes and, as Penrose's nonlinear graviton, for spacetimes with
self-dual or anti-self-dual Weyl curvature; for general real
spacetimes it does not. Local twistors exist on every spacetime.

\textbf{Massive matter.} One twistor describes a massless particle of
helicity \(s\). Two twistors give a timelike momentum
\(p=\pi_1\bar\pi_1+\pi_2\bar\pi_2\), with mass and spin from the
\(SU(2)\) that mixes them; the mass requires a choice of infinity.
Does a two-twistor sector of \(\G^+\) describe massive matter coupled
to the conformal connection?

\textbf{Petrov type.} Can the Petrov type be read from the
eigenvectors of the twistor holonomy?

% --------------------------------------------------
\section{Mathematical Structure}
% --------------------------------------------------

The theory can be summarized by the following objects.

\begin{center}
\renewcommand{\arraystretch}{1.25}
\begin{tabular}{>{\raggedright\arraybackslash}p{0.24\textwidth}
                >{\centering\arraybackslash}p{0.40\textwidth}
                >{\raggedright\arraybackslash}p{0.28\textwidth}}
    \toprule
    Object & Mathematical representation & Purpose\\
    \midrule
    Algebra & \(\Cl(2,4;\R)\) & Underlying conformal geometry\\
    Gauge group & \(\Spin(2,4)\cong SU(2,2)\)
        & Finite conformal transformations\\
    Gauge algebra & \(\so(2,4)\) & Infinitesimal conformal
        transformations\\
    Gauge connection & \(\Om(a,x)\in\Cl^2(2,4)\) & Infinitesimal gauge
        field\\
    Gauge transformation & \(V(x)\in\Spin(2,4)\) & Finite gauge
        transformation\\
    Covariant derivative & \(\D_a\) & Gauge-covariant
        differentiation\\
    Curvature & \(\F(a,b)\in\Cl^2(2,4)\) & Gauge-field strength\\
    Parallel transport & \(T(x_1,x_2)\) & Finite transport between
        points\\
    Coframe & translation sector \(e\wedgeprod\ninf\) & Metric
        \(g=e\cdot e\)\\
    Normal connection & \(f=P\) (Schouten), \(b=0\) & Weyl and Cotton
        curvature\\
    Compensator & \(\mathcal I_\sigma=\sigma\ninf+\nabla\sigma+\rho\nzero\)
        & Choice of scale; Einstein--Hilbert density\\
    Twistor space & \(\T=\G^+P_{\T}\) & Spinors of \(SU(2,2)\); light
        rays\\
    \bottomrule
\end{tabular}
\end{center}

The central relationships are
\[
    \boxed{
        \Om(a,x)
        =
        \frac{1}{2}
        \omega^{AB}(a,x)\,
        \e_A\wedgeprod\e_B
    }
    \qquad
    \boxed{
        V\rev{V}=1
    }
\]
\[
    \boxed{
        \Om'(a)=V\Om(a)\rev V-2(a\cdot\nabla V)\rev V
    }
    \qquad
    \boxed{
        \D_a\Psi
        =
        a\cdot\nabla\Psi
        +
        \frac{1}{2}\Om(a,x)\Psi
    }
\]
\[
    \boxed{
        \F(a,b)
        =
        a\cdot\nabla\Om(b)
        -
        b\cdot\nabla\Om(a)
        +
        \frac{1}{2}[\Om(a),\Om(b)]
    }
    \qquad
    \boxed{
        T=\mathcal P\exp\Bigl(-\tfrac12\int_\gamma\Om\Bigr)
    }
\]
and, for the dynamics,
\[
    \boxed{
        \Om_{\mathrm{normal}}=\Om_M+e\wedgeprod\ninf+P\wedgeprod\nzero
    }
\]
\[
    \boxed{
        S_1+\beta S_2+\gamma S_3
        \ \Longrightarrow\
        G_{\mu\nu}=\Lambda g_{\mu\nu},
        \quad
        \Lambda=6\bigl(\gamma-\tfrac1\beta\bigr),
        \quad
        f=\tfrac1\beta\,e .
    }
\]

% --------------------------------------------------
\section{Conclusion}
% --------------------------------------------------

The formulation extends the algebraic setting of conventional Gauge
Theory Gravity from STA to CSTA. Its principal feature is the
representation of the complete conformal gauge connection as a single
bivector-valued field,
\[
    \boxed{
        \Om(a,x)\in\Cl^2(2,4;\R),
    }
\]
containing the generators of Lorentz transformations, translations,
dilations and special conformal transformations, with finite gauge
transformations and parallel transport represented by normalized CSTA
versors,
\[
    \boxed{
        V(x)\rev{V}(x)=1.
    }
\]

The translation and Lorentz sectors hold GTG's position and rotation
gauge fields, exactly so for Kerr, and GTG's gauge freedom acts as
CSTA gauge transformations. The normal conditions fix the
special-conformal sector to the Schouten tensor, Weyl rescalings are
versor gauge transformations, and the dilation and special-conformal
sectors carry a choice of scale. With the infinity twistor as
compensator, Einstein gravity with a cosmological constant follows in
both second- and first-order form; the action quadratic in curvature
gives Weyl gravity instead, and mixing it in adds a massive spin-2
ghost. The real algebra supplies its own complex structure, and with
it twistor space and Penrose's local twistor transport.

The displacement and rotational gauge structures of conventional GTG
thus form a constrained sector of the conformal formulation.

\begin{thebibliography}{9}
\bibitem{LDG} A.~Lasenby, C.~Doran and S.~Gull, \emph{Gravity, gauge
theories and geometric algebra}, Phil.\ Trans.\ R.\ Soc.\ Lond.\ A
\textbf{356} (1998) 487--582, arXiv:gr-qc/0405033.
\bibitem{CLDG} A.~Challinor, A.~Lasenby, C.~Doran and S.~Gull,
\emph{Massive, non-ghost solutions for the Dirac field coupled
self-consistently to gravity}, Gen.\ Rel.\ Grav.\ \textbf{29} (1997)
1527.
\bibitem{Leong} B.~Leong, \emph{Spin-torsion, braneworlds and changing
symmetry in the universe}, PhD thesis, University of Cambridge (2002),
arXiv:gr-qc/0311016.
\end{thebibliography}

\end{document}
